\chapter{Triển khai}
\label{chap:implementation}

Chương này trình bày thiết kế và triển khai hệ thống streaming Big Data
cho dữ liệu cảm biến IoT. Nội dung được tổ chức theo bốn phần: (i) kiến
trúc tổng thể và sơ đồ thành phần; (ii) data contract và mô hình Bronze--Silver--Gold;
(iii) môi trường và công cụ triển khai; (iv) chi tiết các Spark job và
Airflow DAG.

\section{Kiến trúc tổng thể}

Hệ thống được tổ chức thành hai mức vận hành: \textbf{mức baseline} đảm bảo
pipeline end-to--end chạy được, và \textbf{mức nâng cao} bổ sung các khả năng
đặc trưng của Big Data Systems như partition, checkpoint, fault tolerance và
observability. Về mặt chức năng, kiến trúc mức nâng cao được chia thành
\textbf{năm tầng xử lý} xếp theo chiều dọc, cùng một \textbf{dải
observability} nằm bên phải, xuyên suốt và quan sát tất cả các tầng.
Hình~\ref{fig:architecture} minh hoạ kiến trúc này.

\begin{figure}[H]
\centering
\begin{tikzpicture}[
    box/.style={rectangle, draw, rounded corners=2pt, minimum height=0.8cm,
                align=center, fill=blue!5},
    storage/.style={rectangle, draw, rounded corners=2pt, minimum height=0.8cm,
                    align=center, fill=orange!15},
    orch/.style={rectangle, draw, rounded corners=2pt, minimum height=0.8cm,
                 align=center, fill=purple!10},
    obs/.style={rectangle, draw, rounded corners=2pt, minimum height=0.75cm,
                align=center, fill=pink!10},
    layer/.style={rectangle, draw, dashed, rounded corners=3pt, inner sep=5pt},
    lbl/.style={font=\footnotesize},
    arrow/.style={-{Stealth[length=2.5mm]}, thick},
    trig/.style={-{Stealth[length=2.2mm]}, thick, dashed},
    metr/.style={-{Stealth[length=2.0mm]}, thin, densely dotted},
    every node/.style={font=\small}
]

% ===== Layer 1: Ingestion =====
\node[box, minimum width=2.4cm] (sim) at (-2.6,0.6) {Sensor\\Simulator};
\node[box, minimum width=2.4cm] (external) at (0.4,0.6) {External\\Producers};
\node[box, minimum width=6.0cm] (kafka) at (-1.1,-1.5)
     {Kafka \texttt{sensor.raw}\\(partitions: 3/6/12)};
\begin{pgfonlayer}{background}
  \node[layer, fill=blue!5, fit=(sim)(external)(kafka),
        label=above:{\footnotesize\textbf{T1 --- Ingestion}}] (L1) {};
\end{pgfonlayer}

% ===== Layer 2: Stream processing =====
\node[box, minimum width=7.6cm] (bronzejob) at (-1.1,-4.3)
     {Spark Structured Streaming --- Bronze job\\[1pt]
      {\footnotesize parse JSON $\cdot$ watermark (baseline) \texttt{event\_time} 10 phút}\\
      {\footnotesize dedup theo \texttt{event\_id} $\cdot$ checkpoint trên HDFS}};
\begin{pgfonlayer}{background}
  \node[layer, fill=blue!5, fit=(bronzejob),
        label=above:{\footnotesize\textbf{T2 --- Stream processing}}] (L2) {};
\end{pgfonlayer}

% ===== Layer 3: Storage / Lakehouse =====
\node[storage, minimum width=2.2cm] (bronze) at (-4.0,-7.1) {Bronze\\(Delta)};
\node[storage, minimum width=2.2cm] (silver) at (-0.8,-7.1) {Silver\\(Delta)};
\node[storage, minimum width=2.2cm] (gold) at (2.4,-7.1) {Gold\\(Delta)};
\node[storage, minimum width=3.0cm] (quarantine) at (-2.05,-8.8)
     {\texttt{silver\_quarantine}\\(Delta)};
\begin{pgfonlayer}{background}
  \node[layer, fill=orange!15, fit=(bronze)(silver)(gold)(quarantine),
        label=above:{\footnotesize\textbf{T3 --- Storage / Lakehouse}},
        label=below:{\footnotesize\itshape trên HDFS}] (L3) {};
\end{pgfonlayer}

% ===== Layer 4: Batch & orchestration =====
\node[box, minimum width=1.9cm] (silverjob) at (-4.2,-11.4)
     {Silver job\\[-1pt]{\footnotesize(Airflow trigger)}};
\node[box, minimum width=1.9cm] (goldjob) at (-1.6,-11.4)
     {Gold job\\[-1pt]{\footnotesize(Airflow trigger)}};
\node[box, minimum width=2.2cm] (maint) at (0.9,-11.4)
     {Compaction /\\Retention\\[-1pt]{\footnotesize(Airflow schedule)}};
\node[orch, minimum width=2.2cm] (airflow) at (4.7,-11.4) {Airflow\\(DAG)};
\begin{pgfonlayer}{background}
  \node[layer, fit=(silverjob)(goldjob)(maint)(airflow),
        label={[xshift=1.2cm]above:{\footnotesize\textbf{T4 --- Batch \& orchestration}}}] (L4) {};
\end{pgfonlayer}

% ===== Layer 5: Serving & UI =====
\node[box, minimum width=2.6cm] (sql) at (2.4,-13.7) {Spark SQL\\serving};
\node[box, dashed, fill=purple!5, minimum width=2.6cm] (webui) at (-0.8,-13.7)
     {Web UI\\[-1pt]{\footnotesize(optional)}};
\begin{pgfonlayer}{background}
  \node[layer, fit=(sql)(webui),
        label=above:{\footnotesize\textbf{T5 --- Serving \& UI}}] (L5) {};
\end{pgfonlayer}

% ===== Observability strip (right) =====
\node[obs, minimum width=2.2cm] (prom) at (7.8,-4.3) {Prometheus};
\node[obs, minimum width=2.2cm] (grafana) at (7.8,-11.4) {Grafana};
\node[lbl, align=center] (jmxnote) at (7.8,-12.6) {JMX Exporter\\(Spark \& Kafka)};
\begin{pgfonlayer}{background}
  \node[layer, fill=pink!10, fit=(prom)(grafana)(jmxnote),
        label=above:{\footnotesize\textbf{Observability}}] (LOBS) {};
\end{pgfonlayer}

% ===== Data-flow arrows =====
\draw[arrow] (sim.south) -- (kafka.north);
\draw[arrow] (external.south) -- (kafka.north);
\draw[arrow] ([xshift=2.7cm]kafka.south) --
      node[right=2pt, lbl] {consume} ([xshift=2.7cm]bronzejob.north);
\draw[arrow] (bronzejob.south -| bronze.north) --
      node[left=2pt, lbl] {Bronze job} (bronze.north);
\draw[arrow] (bronze.east) -- node[above, lbl] {Silver job} (silver.west);
\draw[arrow] (silver.east) -- node[below, lbl] {Gold job} (gold.west);
\draw[arrow] (silverjob.north) -- (-4.2,-10.05) --
      node[above, lbl] {record lỗi} (-1.55,-10.05) -- (-1.55,-9.2);
\draw[arrow] (gold.east) -- node[above, lbl] {truy vấn} (6.2,-7.1) --
      (6.2,-13.7) -- (sql.east);

% ===== Orchestration arrow =====
\draw[trig] (airflow.west) -- node[above, lbl] {schedule} (maint.east);

% ===== Metric arrows =====
\draw[metr] (kafka.east) -- node[above, lbl, pos=0.8] {JMX} (prom.north west);
\draw[metr] (bronzejob.east) -- node[above, lbl] {metrics} (prom.west);
\draw[arrow] (prom.south) -- (grafana.north);

% ===== Legend =====
\node[lbl] at (0,-15.0)
     {nét liền: luồng dữ liệu $\cdot$ nét đứt: điều phối (trigger/schedule)
      $\cdot$ nét chấm: metric};
\end{tikzpicture}
\caption{Kiến trúc hệ thống theo năm tầng: T1 Ingestion (Sensor Simulator và
external producer $\rightarrow$ Kafka \texttt{sensor.raw}); T2 Stream processing
(Bronze job: parse JSON, watermark baseline 10 phút, dedup \texttt{event\_id},
checkpoint trên HDFS); T3 Storage/Lakehouse (Bronze/Silver/Gold trên HDFS, nhánh rẽ
\texttt{silver\_quarantine} từ Silver job); T4 Batch \& orchestration (Silver
job và Gold job được Airflow khởi động/dừng theo lịch, compaction/retention do
Airflow schedule); T5 Serving \& UI (Spark SQL và Web UI tuỳ chọn, viền nét
đứt). Dải Observability (Prometheus/Grafana/JMX Exporter) bên phải, quan sát
các tầng. Dữ liệu giữa các tầng được chuyển bởi job ghi nhãn trên mũi tên,
không ``tự chảy''.}
\label{fig:architecture}
\end{figure}

\textbf{Tầng 1 --- Ingestion.} Sensor Simulator sinh tải có kiểm soát với tốc
độ cấu hình được và chế độ giả lập lỗi; các external producer có thể được bổ
sung để mô phỏng nguồn dữ liệu thật. Toàn bộ event được gửi vào Kafka topic
\texttt{sensor.raw} (3/6/12 partition tuỳ kịch bản), đóng vai trò bộ đệm phân
tán, tách khâu sinh dữ liệu khỏi khâu xử lý.

\textbf{Tầng 2 --- Stream processing.} Bronze job là một Spark Structured
Streaming query đọc từ \texttt{sensor.raw}, parse JSON theo schema đã định,
gắn mốc platform receive time thành trường provenance
\texttt{ingest\_time\_spark} (Ct. thời gian $t^{s}$ ở
Mục~\ref{subsec:event-tuple}), áp dụng watermark baseline 10 phút trên
\texttt{event\_time} và loại trùng theo \texttt{event\_id} trong watermark
bound (dedup pass 1). Kết quả được ghi xuống bảng Bronze; trạng thái và offset
của query được lưu checkpoint trên HDFS để khôi phục được sau khi restart.
Giá trị 10 phút là cấu hình baseline cố định; điểm mà một extension thích nghi
có thể can thiệp được mô tả ở Mục~\ref{subsec:bronze-baseline}.

\textbf{Tầng 3 --- Storage / Lakehouse.} HDFS cung cấp lớp lưu trữ phân tán
với replication; trên đó, Delta Lake tổ chức dữ liệu theo ba lớp
Bronze/Silver/Gold. Bản ghi lỗi được tách riêng vào bảng
\texttt{silver\_quarantine} (vật lý: \texttt{delta/silver/sensor\_quarantine};
sẽ mô tả ở mục sau). Mỗi lần ghi là một
transaction ACID trên Delta log, phục vụ cả batch lẫn streaming.

\textbf{Tầng 4 --- Batch \& orchestration.} Silver job đọc Bronze, kiểm tra
ràng buộc giá trị và tách bản ghi lỗi sang \texttt{silver\_quarantine}; Gold
job đọc Silver và tổng hợp theo tumbling window 1 phút/5 phút; các tác vụ
compaction (\texttt{OPTIMIZE}) và retention (\texttt{VACUUM}) gom nhỏ tệp và
dọn dữ liệu cũ. Về mô hình thực thi, Silver job là batch job: đọc bảng Bronze
bằng Spark batch, kiểm tra chất lượng và ghi append vào Silver
(Listing~\ref{lst:silver}), được Airflow chạy theo lịch (task
\texttt{silver\_batch\_validation}). Gold job là streaming query:
\texttt{readStream} từ Silver, tổng hợp theo cửa sổ thời gian rồi ghi append
vào Gold (Mục~\ref{subsec:goldjob}), nhưng cũng được Airflow khởi động và
dừng theo lịch (task \texttt{gold\_refresh}) nên mỗi lần chạy vẫn có điểm bắt
đầu/kết thúc rõ ràng. Riêng Bronze job là streaming query long-running, chạy
độc lập ngoài sự điều phối của Airflow.

\textbf{Tầng 5 --- Serving \& UI.} Kết quả phân tích được đọc trực tiếp từ
bảng Gold qua Spark SQL; một Web UI (tuỳ chọn, không thuộc phạm vi bắt buộc)
có thể được bổ sung để data scientist xem lịch sử event, thống kê theo cửa sổ
và bảng quarantine. Web UI chỉ đọc Gold, không đọc trực tiếp Kafka hay
Bronze.

\textbf{Observability --- dải bên phải.} Xuyên suốt các tầng, Spark và Kafka
expose metric JVM/JMX qua JMX Exporter; tiến độ streaming query của Spark
được scrape qua endpoint tuỳ biến. Prometheus thu thập toàn bộ metric và
hiển thị trên Grafana kèm alert.

Lưu ý rằng dữ liệu giữa các tầng không ``tự chảy'': mỗi mũi tên trên
Hình~\ref{fig:architecture} tương ứng một job cụ thể --- Bronze job (streaming,
chạy độc lập), Silver job (batch job do Airflow chạy theo lịch) và Gold job
(streaming query được Airflow khởi động/dừng theo lịch).
Hình~\ref{fig:event-flow} theo dõi chi tiết hành trình của một event đơn lẻ
qua kiến trúc này.

\begin{figure}[H]
\centering
\begin{tikzpicture}[
    box/.style={rectangle, draw, rounded corners=2pt, minimum height=0.75cm,
                minimum width=4.2cm, align=center, fill=blue!5},
    storage/.style={rectangle, draw, rounded corners=2pt, minimum height=0.75cm,
                    minimum width=4.2cm, align=center, fill=orange!15},
    bad/.style={rectangle, draw, rounded corners=2pt, minimum height=0.75cm,
              minimum width=3.4cm, align=center, fill=red!8},
    serving/.style={rectangle, draw, rounded corners=2pt, minimum height=0.75cm,
                    minimum width=4.2cm, align=center, fill=purple!5},
    note/.style={align=left, font=\footnotesize\itshape},
    arrow/.style={-{Stealth[length=2.5mm]}, thick},
    lbl/.style={font=\footnotesize},
    every node/.style={font=\small}
]

% Main vertical chain
\node[box] (sensor) at (0,0) {Sensor (producer)};
\node[note, anchor=west] at (2.5,0)
     {phát sinh event, gắn \texttt{event\_time}\\
      (thời điểm thực tế tại cảm biến)};
\node[box] (kafka) at (0,-1.7) {Kafka \texttt{sensor.raw}};
\node[box] (bjob) at (0,-3.4) {Bronze job};
\node[note, anchor=west] at (2.5,-3.4)
     {gắn \texttt{ingest\_time\_spark} lúc Spark nhận\\
      watermark baseline \texttt{event\_time}: 10 phút\\
      dedup theo \texttt{event\_id} trong bound};
\node[storage] (bronze) at (0,-5.1) {Bảng Bronze (Delta)};
\node[box] (sjob) at (0,-6.8) {Silver job (validate)};
\node[storage] (silver) at (0,-8.5) {Bảng Silver (bản ghi hợp lệ)};
\node[box] (gjob) at (0,-10.2) {Gold job (tumbling window 1 phút / 5 phút)};
\node[storage] (gold) at (0,-11.9) {Bảng Gold (aggregate)};
\node[serving] (serve) at (0,-13.6) {Spark SQL / Web UI \footnotesize(optional)};

% Quarantine branch
\node[bad] (quarantine) at (-5.2,-6.8) {\texttt{silver\_quarantine}};

% Arrows
\draw[arrow] (sensor) -- (kafka);
\draw[arrow] (kafka) -- (bjob);
\draw[arrow] (bjob) -- (bronze);
\draw[arrow] (bronze) -- (sjob);
\draw[arrow] (sjob) -- node[right=2pt, lbl] {hợp lệ} (silver);
\draw[arrow] (sjob.west) -- node[above, lbl] {record lỗi} (quarantine.east);
\draw[arrow] (silver) -- (gjob);
\draw[arrow] (gjob) -- (gold);
\draw[arrow] (gold) -- (serve);
\end{tikzpicture}
\caption{Luồng sự kiện: hành trình của một event từ cảm biến qua các tầng.
\texttt{event\_time} phát sinh tại sensor; \texttt{ingest\_time} do
producer/gateway gắn trước khi gửi, và \texttt{ingest\_time\_spark} được
Bronze job gắn khi nhận event từ Kafka (cơ sở tính ingestion lateness;
xem Bảng~\ref{tab:time-taxonomy}). Watermark baseline 10 phút giới hạn cửa sổ
dedup theo \texttt{event\_id}; bản ghi vi phạm quy tắc chất lượng bị Silver
job tách sang \texttt{silver\_quarantine} thay vì vào Silver.}
\label{fig:event-flow}
\end{figure}

\subsection{Các thành phần chính và vai trò}

\begin{table}[H]
\centering
\caption{Vai trò và lý do lựa chọn của từng thành phần trong pipeline.}
\label{tab:components}
\begin{tabularx}{\textwidth}{l l X}
\toprule
\textbf{Thành phần} & \textbf{Vai trò} & \textbf{Lý do lựa chọn} \\
\midrule
Sensor Simulator & Sinh tải có kiểm soát & Chủ động tạo 1k--50k event/s \\
Kafka & Distributed ingestion/buffer & Tách producer \& processor, partition để scale \\
Spark SS & Distributed stream computation & Window, watermark, dedup, stateful \\
HDFS & Distributed persistent storage & Replication, chịu lỗi, persistent checkpoint \\
Delta Lake & Lakehouse layer & ACID, schema, batch + stream, time travel \\
Airflow & Batch orchestration & Khởi động/dừng Silver \& Gold job theo lịch; compaction, retention, quality check \\
Prometheus & Time-series metric store & Hỗ trợ pull model, tích hợp JMX \\
JMX Exporter & JVM metrics $\rightarrow$ Prometheus & Kafka chủ yếu expose JMX \\
Grafana & Dashboard + alert & Quan sát throughput, lag, latency \\
Spark SQL & Analytical serving & Truy vấn trực tiếp Delta table \\
Web UI (optional) & Dashboard cho data scientist & Streamlit/FastAPI, đọc Gold; không bắt buộc triển khai \\
\bottomrule
\end{tabularx}
\end{table}

\subsection{Tầng serving và Web UI (tuỳ chọn)}

Tầng serving là điểm kết thúc của pipeline dữ liệu: mọi truy vấn phân tích
đều xuất phát từ bảng Gold --- nơi đã chứa aggregate theo cửa sổ 1 phút và
5 phút, sạch trùng và đã qua kiểm tra chất lượng. Data scientist có thể dùng
Spark SQL trực tiếp để xem lịch sử event, thống kê theo
\texttt{sensor\_id}/\texttt{location}, hoặc đối chiếu tỷ lệ bản ghi trong
\texttt{silver\_quarantine}.

Như một phần mở rộng \emph{tuỳ chọn}, một Web UI có thể được triển khai để
đóng gói các truy vấn Spark SQL thành dashboard thân thiện hơn. Công nghệ
gợi ý là Streamlit (làm dashboard nhanh, thuần Python) hoặc FastAPI + một
frontend nhẹ (khi cần API riêng cho nhiều client). Nguyên tắc thiết kế là
Web UI chỉ đọc bảng Gold qua Spark SQL --- không đọc trực tiếp Kafka hay
Bronze --- để giữ nhất quán một nguồn dữ liệu đã kiểm chứng và không tạo
thêm tải lên tầng ingestion. Web UI không thuộc phạm vi bắt buộc của đồ án;
giá trị chính của đề tài nằm ở năm tầng xử lý và phương pháp benchmark.

\section{Data contract và mô hình Bronze--Silver--Gold}

\subsection{Sensor event schema}

Schema sự kiện được thiết kế tối giản nhưng đủ trường để phục vụ các
RQ về dedup, out-of-order, latency và data quality. Contract gồm 9 trường
khai báo như Listing~\ref{lst:schema}; riêng Bronze job gắn thêm trường
provenance \texttt{ingest\_time\_spark} (thời điểm nền tảng nhận,
Bảng~\ref{tab:time-taxonomy}) ngoài contract này:

\begin{lstlisting}[language=yaml,basicstyle=\ttfamily\footnotesize,frame=single,
  caption={Sensor event schema (JSON).},label={lst:schema}]
{
  "event_id":     "9c0e12ab-3f4d-4a8b-b1c6-7e2f0a9d4c11",
  "sensor_id":    "sensor-001923",
  "event_time":   "2026-08-24T10:15:32.120Z",
  "ingest_time":  "2026-08-24T10:15:32.202Z",
  "sensor_type":  "temperature",
  "value":        31.42,
  "unit":         "C",
  "location":     "zone-04",
  "sequence_no":  91283
}
\end{lstlisting}

Trong đó: \texttt{event\_id} phục vụ deduplication; \texttt{event\_time} là
mốc thời gian thực tế tại cảm biến (dùng cho event-time window); \texttt{ingest\_time}
là mốc producer/gateway gắn trước khi gửi (dùng tách phần trễ ngoài nền tảng
khỏi ingestion lateness, xem Bảng~\ref{tab:time-taxonomy});
\texttt{sequence\_no} giúp phát hiện mất event hoặc out-of-order.

\subsection{Mô hình Bronze--Silver--Gold}

Dữ liệu được tổ chức thành ba lớp trên HDFS, mỗi lớp là một bảng Delta:

\begin{table}[H]
\centering
\caption{Đặc tả ba lớp dữ liệu Bronze/Silver/Gold.}
\label{tab:medallion}
\begin{tabularx}{\textwidth}{l l X}
\toprule
\textbf{Lớp} & \textbf{Schema} & \textbf{Nội dung \& quy tắc} \\
\midrule
Bronze & Raw payload + metadata &
  Mọi event hợp lệ về mặt JSON, kèm \texttt{ingest\_time}, partition theo \texttt{event\_date}. \\
Silver & Validated + deduplicated &
  Schema enforced; dedup pass 2 theo \texttt{event\_id} trong phạm vi batch
  (không phụ thuộc watermark); record vi phạm required-field hoặc physical
  bounds bị tách vào bảng \texttt{silver\_quarantine} kèm lý do. \\
Gold & Window aggregates &
  avg / min / max / count của \texttt{value} theo 1-minute và 5-minute
  tumbling window, nhóm theo \texttt{location} $\times$ \texttt{sensor\_type},
  kèm \texttt{window\_start}/\texttt{window\_end}; partition theo
  \texttt{event\_date}. \\
\bottomrule
\end{tabularx}
\end{table}

\textbf{Quy tắc phân vùng:} Delta Lake khuyến nghị tránh partition có
cardinality quá cao~\cite{delta_lake_docs}. Việc partition theo
\texttt{sensor\_id} với hàng trăm nghìn sensor sẽ tạo ra vô số thư mục
nhỏ, không khả thi. Hệ thống đề xuất partition theo \texttt{event\_date}
hoặc \texttt{event\_date/hour} khi data đủ lớn, đảm bảo mỗi partition
có kích thước hợp lý (mức khuyến nghị thực dụng $\sim$1 GB~\cite{delta_lake_docs}).

\section{Môi trường và công cụ}

\subsection{Phiên bản phần mềm}

\begin{table}[H]
\centering
\caption{Phiên bản phần mềm sử dụng trong triển khai.}
\label{tab:versions}
\begin{tabularx}{\textwidth}{l l l X}
\toprule
\textbf{Thành phần} & \textbf{Phiên bản} & \textbf{Image} & \textbf{Ghi chú} \\
\midrule
Apache Kafka & 4.3.x & \texttt{apache/kafka:4.3.1} & KRaft mode, không cần ZooKeeper~\cite{kafka_docs_4_3} \\
Apache Spark & 4.2.x & custom Spark 4.2 image & Structured Streaming + Delta connector \\
Delta Lake & tương thích Spark 4.2 & jar trong Spark image & ACID + streaming integration~\cite{delta_lake_docs} \\
Hadoop HDFS & 3.5.x & \texttt{apache/hadoop:3.5} & NameNode + DataNode~\cite{hdfs_docs_3_5} \\
Apache Airflow & 2.x & official Docker image & Batch/maintenance orchestration~\cite{airflow_docs} \\
Prometheus & mới nhất & official image & Pull model, scrape mỗi 15s \\
Grafana & mới nhất & official image & Dashboard + alert~\cite{grafana_alerting} \\
JMX Exporter & mới nhất & \texttt{prometheus/jmx\_exporter} & Bridge JVM $\rightarrow$ Prometheus~\cite{prometheus_jmx_exporter} \\
\bottomrule
\end{tabularx}
\end{table}

\subsection{Cấu hình Docker Compose}

Hệ thống được đóng gói bằng một file \texttt{docker-compose.yml} với
các logical service sau:

\begin{lstlisting}[language=yaml,basicstyle=\ttfamily\footnotesize,frame=single,
  caption={Logical services trong docker-compose.yml.},label={lst:compose}]
# --- Data plane ---
sensor-simulator:        # Python producer; gui event vao sensor.raw
kafka-1:                 # KRaft single broker (phase A), 3 brokers (phase B)
spark-master:            # Spark driver + master UI
spark-worker-1:          # Spark executor
spark-worker-2:          # Spark executor

# --- Storage ---
hdfs-namenode:           # NameNode, expose 9870 (UI), 9000 (RPC)
hdfs-datanode-1:         # DataNode, replication factor = 2
hdfs-datanode-2:         # DataNode

# --- Orchestration ---
airflow-webserver:       # Airflow UI
airflow-scheduler:       # Airflow scheduler
postgres-airflow:        # Airflow metadata DB

# --- Observability ---
prometheus:              # Scrape config: kafka, spark, jmx-exporter
grafana:                 # Provisioning Prometheus datasource
jmx-kafka:               # JMX exporter sidecar cho kafka-1
kafka-exporter:          # Consumer lag metric
spark-metrics:           # Custom endpoint expose Spark query progress
\end{lstlisting}

\textbf{Mạng và persistent volume:} toàn bộ service giao tiếp qua một
\texttt{bridge network} nội bộ. Hai volume quan trọng được mount:
\texttt{hdfs\_data/} (lưu Delta tables và HDFS namenode/datanode metadata)
và \texttt{spark\_checkpoints/} (lưu checkpoint location, đặt tên bắt
đầu bằng dấu gạch dưới \texttt{\_checkpoints/} để Delta VACUUM không
coi chúng là data file)~\cite{delta_lake_docs}.

\section{Cài đặt chi tiết}

\subsection{Sensor producer}

Producer được viết bằng Python sử dụng \texttt{kafka-python} hoặc
\texttt{confluent-kafka-python}, gửi message dạng JSON với cấu hình:

\begin{lstlisting}[language=yaml,basicstyle=\ttfamily\footnotesize,frame=single,
  caption={Cấu hình producer.},label={lst:producer}]
acks: all                         # wait for all in-sync replicas
enable.idempotence: true          # producer avoids duplicate sequence
compression.type: lz4
linger.ms: 5                      # batch before send
batch.size: 65536                 # 64 KB
max.in.flight.requests.per.connection: 5
\end{lstlisting}

Hai chế độ vận hành: \texttt{--rate N} để giữ tốc độ ổn định $N$ event/s,
\texttt{--fault inject} để cố tình đưa vào các tỷ lệ bad record theo
một profile định trước (ví dụ 95\% valid, 2\% duplicate, 1\% missing field,
1\% temperature = 9999, 0.5\% late 30s, 0.5\% late 5min).

\subsection{Spark streaming job (Bronze)}
\label{subsec:bronze-baseline}

Job Bronze đọc từ Kafka, parse JSON, áp dụng schema, watermark và
deduplication, ghi xuống bảng Bronze trên HDFS. Pseudo-code (dạng rút gọn
để trình bày; bản cài đặt thực tế gắn thêm trường provenance
\texttt{ingest\_time\_spark} bằng \texttt{current\_timestamp()} trước khi ghi):

\begin{lstlisting}[language=yaml,basicstyle=\ttfamily\footnotesize,frame=single,
  caption={Pseudo-code cho Spark Bronze job.},label={lst:bronze}]
raw = (spark.readStream
         .format("kafka")
         .option("kafka.bootstrap.servers", "kafka-1:9092")
         .option("subscribe", "sensor.raw")
          .option("startingOffsets", "earliest")
         .option("failOnDataLoss", "false")
         .load())

parsed = (raw
            .selectExpr("CAST(value AS STRING) AS json_str")
            .select(from_json("json_str", sensor_schema).alias("e"))
            .select("e.*")
            .withColumn("event_date", to_date("event_time")))

validated = (parsed
               .withWatermark("event_time", "10 minutes")
               .dropDuplicates(["event_id"]))   # dedup inside watermark bound

(validated.writeStream
   .format("delta")
   .mode("append")
   .option("checkpointLocation", "hdfs://nn:9000/_checkpoints/bronze")
   .partitionBy("event_date")
   .trigger(processingTime="5 seconds")  # long-running baseline
   # .trigger(availableNow=True)         # drain mode for scheduled refresh
   .start("hdfs://nn:9000/delta/bronze/sensor"))
\end{lstlisting}

Điểm cần lưu ý:
\begin{itemize}
  \item \texttt{withWatermark} kết hợp với \texttt{dropDuplicates} cho
        phép dedup có giới hạn thời gian, tránh giữ state vô hạn~\cite{spark_docs_4_2}.
        Giá trị 10 phút là \textbf{cấu hình baseline cố định}; ngữ nghĩa
        watermark (tiến độ query-level, đơn điệu, không tự xoá event muộn)
        theo Mục~\ref{subsec:watermark};
  \item \textbf{Cột cơ sở watermark.} Bản cài đặt thực tế tạo cột phụ
        \texttt{event\_ts} = \texttt{event\_time} ép kiểu timestamp, với
        fallback sang \texttt{kafka\_ts} (timestamp do broker gắn) khi
        \texttt{event\_time} rỗng; watermark và dedup chạy trên
        \texttt{event\_ts}, sau đó cột phụ bị loại trước khi ghi. Nhờ
        fallback, bản ghi thiếu \texttt{event\_time} vẫn được nhận diện trùng
        trong bound và vẫn vào Bronze để Silver phân loại về sau;
  \item \texttt{checkpointLocation} trên HDFS đảm bảo checkpoint tồn
        tại qua restart container;
  \item \textbf{Trigger.} Bản cài đặt baseline chạy Bronze job long-running với
        \texttt{trigger(processingTime="5 seconds")}. Variant
        \texttt{trigger(availableNow=True)} --- xử lý hết dữ liệu hiện có rồi
        tự dừng --- được dùng ở chế độ drain/batch refresh, cho phép Airflow
        khởi động và dừng job theo lịch như một job có điểm bắt đầu/kết thúc
        rõ ràng, trong khi mô hình thực thi vẫn là streaming query.
\end{itemize}

\textbf{Điểm can thiệp (seam) của extension thích nghi.} Theo thiết kế chính
sách ở Mục~\ref{subsec:policy}, tham số watermark (và trigger) là cấu hình khởi
tạo của query. Một controller bên ngoài --- không nằm trong bản cài đặt hiện
tại --- có thể đọc metric chất lượng từ dải observability, chọn lại tham số và
khởi động lại query giữa các epoch. Seam này không thay đổi mô hình thực thi:
Bronze job vẫn là streaming query, watermark vẫn là watermark query-level của
Spark, và extension là policy epoch chứ không phải per-record adaptive watermark.
Khi đổi tham số giữa các epoch, controller phải tuân theo giao thức liên tục
được đề xuất ở Mục~\ref{subsec:policy} (không hạ frontier watermark hiệu
dụng); restart một query không tự động bảo đảm điều đó.

\subsection{Spark job (Bronze $\rightarrow$ Silver)}

Silver job là một \textbf{batch job} do Airflow khởi chạy theo lịch: đọc bảng
Bronze bằng Spark batch, áp dụng quy tắc chất lượng (required-field, physical
bounds), ghi bản hợp lệ vào Silver và tách bản ghi vi phạm sang bảng
quarantine kèm lý do. Pass này thực hiện dedup lần hai trên
\texttt{event\_id} trong phạm vi batch hiện tại (bắt các bản trùng đến muộn
qua cửa sổ dedup pass 1) và ghi kết quả theo cách idempotent --- bản ghi chỉ
được thêm khi \texttt{event\_id} chưa tồn tại trong bảng đích --- để các lần
chạy lại không tạo bản sao:

\begin{lstlisting}[language=yaml,basicstyle=\ttfamily\footnotesize,frame=single,
  caption={Pseudo-code cho Bronze $\rightarrow$ Silver job.},label={lst:silver}]
bronze = spark.read.format("delta").load("hdfs://nn:9000/delta/bronze/sensor")

valid = (bronze
           .filter(col("value").between(-50, 200))     # physical bounds
           .filter(col("event_time").isNotNull())
           .dropDuplicates(["event_id"]))              # second-pass dedup

quarantine = (bronze
                .subtract(valid)
                .withColumn("quarantine_reason",
                            concat_ws(",", ...)))

(valid.write.format("delta").mode("append")
       .partitionBy("event_date")
       .save("hdfs://nn:9000/delta/silver/sensor"))

(quarantine.write.format("delta").mode("append")
       .save("hdfs://nn:9000/delta/silver/sensor_quarantine"))
# both writes are idempotent keyed inserts (merge on event_id),
# so re-running the batch pass does not create duplicates
\end{lstlisting}

\subsection{Spark job (Silver $\rightarrow$ Gold)}
\label{subsec:goldjob}

Job Gold thực hiện window aggregation. Trong bản minh hoạ này job được viết
dưới dạng streaming query (đọc/ghi bằng \texttt{readStream}/\texttt{writeStream}
với watermark baseline 10 phút), nhưng được Airflow khởi động và dừng theo lịch
(task \texttt{gold\_refresh} trong Mục~\ref{subsec:airflowdag}) nên mỗi lần
chạy vẫn có điểm bắt đầu/kết thúc rõ ràng như một job batch:

\begin{lstlisting}[language=yaml,basicstyle=\ttfamily\footnotesize,frame=single,
  caption={Pseudo-code cho Silver $\rightarrow$ Gold job.},label={lst:gold}]
silver = (spark.readStream.format("delta")
            .load("hdfs://nn:9000/delta/silver/sensor")
            .withWatermark("event_time", "10 minutes"))

agg_1m = (silver
            .groupBy(window("event_time", "1 minute"),
                     col("location"),
                     col("sensor_type"))
            .agg(avg("value").alias("avg_value"),
                 min("value").alias("min_value"),
                 max("value").alias("max_value"),
                 count("*").alias("event_count")))

(agg_1m.writeStream
   .format("delta")
   .outputMode("append")
   .option("checkpointLocation", "hdfs://nn:9000/_checkpoints/gold_1m")
   .partitionBy("event_date")
   .start("hdfs://nn:9000/delta/gold/sensor_1m"))
\end{lstlisting}

Cùng một pattern lặp lại cho 5-minute window. Đầu ra Gold được query
bằng Spark SQL cho các dashboard. Như đã nói ở Mục~\ref{subsec:goldjob},
tuy cài đặt bằng streaming API, job này được Airflow khởi động và dừng
theo lịch nên mỗi lần chạy có điểm bắt đầu/kết thúc rõ ràng, khác với
Bronze job chạy liên tục.

Về ngữ nghĩa đầu ra: query Gold chạy ở \texttt{outputMode("append")} với
watermark, nên aggregate của mỗi cửa sổ được phát hành đúng một lần khi
watermark vượt qua thời điểm kết thúc cửa sổ. Event đến sau thời điểm đó với
event time thuộc cửa sổ đã chốt bị engine lọc khỏi phép tổng hợp và không làm
thay đổi dòng đã công bố (Mục~\ref{subsec:watermark}); event đó cũng không tự
động vào quarantine. Vì vậy Gold append trong baseline \emph{không} phải cơ
chế sửa kết quả phổ quát: việc phát hành lại giá trị cửa sổ (CORRECTED) là
mục tiêu của extension đề xuất (Mục~\ref{subsec:lifecycle}), không phải hành
vi mặc định của query hiện tại.

\subsection{Airflow DAG}
\label{subsec:airflowdag}

Airflow DAG không can thiệp vào Bronze job (streaming query long-running ở
tầng 2). DAG chỉ quản lý các tác vụ có điểm bắt đầu/kết thúc rõ ràng, bao
gồm việc khởi động/dừng Silver job và Gold job theo lịch:

\begin{lstlisting}[language=yaml,basicstyle=\ttfamily\footnotesize,frame=single,
  caption={Cấu trúc Airflow DAG.},label={lst:airflow}]
iot_lakehouse_maintenance_dag:
  schedule: "@hourly"
  tasks:
    - check_raw_data            # verify Kafka has new events
    - bronze_quality_check      # count bad record ratio in Bronze
    - silver_batch_validation   # check schema and uniqueness
    - gold_refresh              # trigger Silver -> Gold job if needed
    - delta_compaction          # OPTIMIZE on Gold/Silver
    - retention_check           # VACUUM with safe retention horizon
    - benchmark_report          # export hourly metric report
\end{lstlisting}

\textbf{Lưu ý vận hành:} \texttt{VACUUM} có thể xoá các tệp mà streaming
query hoặc job batch vẫn cần nếu retention quá ngắn. Theo tài liệu
Delta~\cite{delta_lake_docs}, nên giữ retention $\geq$ 7 ngày trong giai
đoạn phát triển.

\subsection{Cấu hình quan trọng}

Một số cấu hình then chốt được thống nhất trong toàn pipeline:

\begin{table}[H]
\centering
\caption{Các cấu hình then chốt và lý do.}
\label{tab:configs}
\begin{tabularx}{\textwidth}{l l X}
\toprule
\textbf{Mục} & \textbf{Giá trị} & \textbf{Lý do} \\
\midrule
Kafka \texttt{acks} & all & Đảm bảo không mất event khi broker fail-over~\cite{kafka_docs_4_3} \\
Kafka \texttt{retention.ms} & 7 ngày & An toàn cho replay/debugging \\
Spark trigger interval & Bronze 5s / Gold 10s & Cân bằng giữa latency và micro-batch overhead~\cite{spark_docs_4_2} \\
Spark \texttt{maxOffsetsPerTrigger} & tỷ lệ với rate & Tránh batch quá lớn khi rate tăng đột biến \\
Watermark bound (baseline) & 10 phút & Đủ lớn cho network delay, đủ nhỏ để giải phóng state \\
Delta checkpoint dir & \texttt{\_checkpoints/} & Tránh bị VACUUM coi là data file~\cite{delta_lake_docs} \\
Delta partition & \texttt{event\_date} & Tránh cardinality cao~\cite{delta_lake_docs} \\
HDFS replication & 2 (dev) / 3 (prod) & Trade-off giữa dung lượng và chịu lỗi~\cite{hdfs_docs_3_5} \\
\bottomrule
\end{tabularx}
\end{table}

\section{Tóm tắt chương}

Chương này đã trình bày:
\begin{itemize}
  \item Kiến trúc năm tầng (Ingestion, Stream processing, Storage/Lakehouse,
        Batch \& orchestration, Serving \& UI) kèm dải observability bên phải;
        Web UI là thành phần tuỳ chọn; dữ liệu giữa các tầng được chuyển bởi
        ba job: Bronze job (streaming, chạy độc lập) cùng Silver job và Gold
        job (được Airflow khởi động/dừng theo lịch), cộng các tác vụ
        compaction/retention do Airflow schedule;
  \item Data contract cho sensor event với 9 trường cốt lõi, cộng trường
        provenance \texttt{ingest\_time\_spark} gắn ở Bronze;
  \item Mô hình Bronze--Silver--Gold, partition theo \texttt{event\_date},
        bản ghi lỗi tách vào \texttt{silver\_quarantine};
  \item Môi trường Docker Compose với các phiên bản phần mềm cụ thể;
  \item Pseudo-code cho ba Spark job (Bronze, Silver, Gold) và Airflow DAG.
\end{itemize}

Chương tiếp theo sẽ trình bày phương pháp benchmark và phân tích số liệu
minh họa cho năm RQ đã đặt ra ở Chương~\ref{chap:introduction}.